% NeurIPS 2026 poster: Tokens-per-Parameter Coverage Is Critical for Robust LLM Scaling Law Extrapolation.
% Build: make (from this folder), with xelatex. A0 landscape (118.9 x 84.1 cm), inside the NeurIPS limit of
% 140 x 90 cm. Theme: gemini (theme/, MIT). Logos come from ../images (SVG, converted by the Makefile);
% figures from figures/ (the paper's Figure 1, panels b and c). Every number and formula is the paper's.
\documentclass[final]{beamer}
\usepackage[size=custom,width=118.9,height=84.1,scale=1.16]{beamerposter}
\usetheme{gemini}
\usecolortheme{gemini}
\usepackage{amsmath,amssymb,booktabs,graphicx,tikz}
% Images are found through TEXINPUTS, set by the Makefile (figures/, build/assets/, ../images/).

% Body text in Lato Regular (the theme's Light is too thin for print); headings stay in Raleway.
\setsansfont{Lato}[UprightFont=*-Regular, ItalicFont=*-Italic, BoldFont=*-Bold, BoldItalicFont=*-BoldItalic]
% CO blue and NC red as in the paper's Figure 1 (matplotlib C0 and C3); orange for the headline numbers only.
\definecolor{co}{RGB}{31, 119, 180}
\definecolor{nc}{RGB}{214, 39, 40}
\definecolor{accent}{RGB}{212, 69, 0}
\newcommand{\CO}{\textcolor{co}{\textbf{CO}}}
\newcommand{\NC}{\textcolor{nc}{\textbf{NC}}}
\setbeamercolor{headline}{fg=darkblue,bg=white}
\setbeamercolor{headline title}{fg=darkblue}
\setbeamercolor{headline rule}{bg=darkblue}
\setbeamercolor{block title alerted}{fg=darkblue,bg=lightblue}
\setbeamerfont{headline title}{size=\fontsize{100}{116}\selectfont,series=\bfseries}
\setbeamerfont{headline author}{size=\fontsize{44}{54}\selectfont}
\setbeamerfont{headline institute}{size=\fontsize{40}{48}\selectfont}
\setbeamertemplate{itemize items}[circle]
\setlength{\leftmargini}{1.2em}

% Three columns: 4 x sep + 3 x col = paper width.
\newlength{\sepwidth}
\newlength{\colwidth}
\setlength{\sepwidth}{0.025\paperwidth}
\setlength{\colwidth}{0.3\paperwidth}
\newcommand{\separatorcolumn}{\begin{column}{\sepwidth}\end{column}}
\newcommand{\statement}[1]{{\usebeamercolor[fg]{block title}\bfseries #1.}\ }

\title{Tokens-per-Parameter Coverage Is Critical for\\Robust LLM Scaling Law Extrapolation}
\author{Joshua Shay Kricheli\textsuperscript{1,$\|$} \quad Alexander Lawrence Reid\textsuperscript{1,$\|$} \quad
  Soumajyoti Sarkar\textsuperscript{2,$*$} \\ Venkata Gandikota\textsuperscript{1} \quad Paulo Shakarian\textsuperscript{1}}
\institute{\textsuperscript{1}Syracuse University \qquad \textsuperscript{2}Amazon AGI Foundations}

\footercontent{\href{https://lab-v2.github.io/tpp-coverage/}{lab-v2.github.io/tpp-coverage} \hfill
  \textsuperscript{$\|$}Equal contribution. \quad \textsuperscript{$*$}This work does not relate to the author's position at Amazon. \hfill
  NeurIPS 2026, Sydney, Australia \hfill \href{https://arxiv.org/abs/2605.08541}{arXiv:2605.08541}}

% Header corners: the labelled QR code of the project page and the NeurIPS logo (left); the Syracuse AI
% Institute and the Leibniz Lab (right). The Leibniz PNG (1024 x 348 px) is set 2.9 cm high, so 305 ppi.
\logoleft{\begin{tabular}{@{}c@{}}\includegraphics[height=6.5cm]{qr_project_page.pdf}\\[-0.2ex]{\normalsize Project page}\end{tabular}\hspace{2cm}\raisebox{2.4cm}{\includegraphics[height=3.6cm]{neurips.pdf}}}
\logoright{\begin{tabular}{@{}c@{}}\includegraphics[height=2.9cm]{syracuse_ai_institute.pdf}\\[2ex]\includegraphics[height=2.9cm]{leibniz_lab.png}\end{tabular}}

% The two designs in the (N, D) plane, from the grids of the paper's Appendix F.1 (grid:ncgrid, grid:tpp):
% sizes N (M parameters), CO ratios k, NC token counts D (M tokens). Holdout: the N = 55.5M row, the
% k = 6 (CO) or D = 300.4M (NC) column, and the three largest sizes at k >= 6.2 (CO) or D >= 325.4M (NC).
\newcommand{\designplane}{%
\begin{tikzpicture}[x=0.36cm, y=0.036cm, font=\normalsize]
  \def\trainN{5.0,10.1,14.4,19.4,22.8,28.8,32.5,35.4,40.3,47.9}
  \def\allN{5.0,10.1,14.4,19.4,22.8,28.8,32.5,35.4,40.3,47.9,55.5,63.9,68.1,76.5}
  \def\bigN{63.9,68.1,76.5}
  \def\trainK{1,1.5,1.9,2,2.5,2.7,3,3.3,3.5,4,4.5,5}
  \def\allK{1,1.5,1.9,2,2.5,2.7,3,3.3,3.5,4,4.5,5,6,6.2,6.5,6.7,7}
  \def\testK{6.2,6.5,6.7,7}
  \def\trainD{10.1,35.1,60.2,80.3,100.4,130.3,150.7,180.6,200.8,225.0,250.9,275.9}
  \def\allD{10.1,35.1,60.2,80.3,100.4,130.3,150.7,180.6,200.8,225.0,250.9,275.9,300.4,325.4,350.4,375.8,400.7}
  \def\bigD{325.4,350.4,375.8,400.7}
  \def\sq{+(-5pt,-5pt) rectangle +(5pt,5pt)}
  \def\hsq{+(-7pt,-7pt) rectangle +(7pt,7pt)}
  \begin{scope}
    \clip (0,0) rectangle (80,560);
    \draw[gray!25] (0,0) grid[xstep=20, ystep=100] (80,560);
    \foreach \k in \trainK { \draw[co!40, line width=1.2pt] (0,0) -- (80,{80*\k}); }
  \end{scope}
  \foreach \n in \trainN { \foreach \d in \trainD { \fill[nc] (\n,\d) \sq; } }
  \foreach \n in \trainN { \foreach \k in \trainK { \fill[co] (\n,{\k*\n}) circle[radius=5.5pt]; } }
  \foreach \n in \allN { \draw[co, line width=1.4pt] (\n,{6*\n}) circle[radius=7.5pt];
    \draw[nc, line width=1.4pt] (\n,300.4) \hsq; }
  \foreach \k in \allK { \draw[co, line width=1.4pt] (55.5,{55.5*\k}) circle[radius=7.5pt]; }
  \foreach \d in \allD { \draw[nc, line width=1.4pt] (55.5,\d) \hsq; }
  \foreach \n in \bigN { \foreach \k in \testK { \draw[co, line width=1.4pt] (\n,{\k*\n}) circle[radius=7.5pt]; }
    \foreach \d in \bigD { \draw[nc, line width=1.4pt] (\n,\d) \hsq; } }
  \draw[thick] (0,0) rectangle (80,560);
  \foreach \x in {0,20,40,60,80} { \node[below] at (\x,0) {\x}; }
  \foreach \y in {0,100,200,300,400,500} { \node[left] at (0,\y) {\y}; }
  \node[below=1.3cm] at (40,0) {Parameters $N$ (M)};
  \node[rotate=90] at (-9,280) {Tokens $D$ (M)};
  \begin{scope}[shift={(2.5,525)}]
    \fill[white, opacity=0.85] (-1.5,22) rectangle (44,-112);
    \fill[co] (0,0) circle[radius=5.5pt]; \node[right] at (1.2,0) {\CO{} training: 12 rays $D = kN$, $k = 1$ to $5$};
    \fill[nc] (0,-42) \sq; \node[right] at (1.2,-42) {\NC{} training: $10 \times 12$ grid};
    \draw[black, line width=1.4pt] (0,-84) circle[radius=7.5pt]; \node[right] at (1.2,-84) {held-out runs};
  \end{scope}
\end{tikzpicture}}

\begin{document}
\begin{frame}[t]
\begin{columns}[t]
\separatorcolumn

% ==== column 1
\begin{column}{\colwidth}

\begin{alertblock}{TL;DR}
  Scaling laws fit at fixed tokens-per-parameter (TPP) ratios $D = kN$ are ill-conditioned when the $N$- and
  $D$-exponents are close. Spreading runs over the $(N, D)$ plane fixes it: \NC{} beats \CO{} on held-out
  runs in \textcolor{accent}{\textbf{97.3\%}} of comparisons, \textcolor{accent}{\textbf{98.7\%}} in BF16.
\end{alertblock}

\begin{block}{Two designs in the $(N, D)$ plane}
  \centering
  \designplane\par\medskip
  \raggedright
  \CO{} (collinear): a fan of rays. \NC{} (non-collinear): a grid, its ratio $k = D/N$ free.
\end{block}

\begin{block}{The law and one ray}
  \[ L(N, D) = A N^{-\alpha} + B D^{-\beta} + E \]
  With $D = kN$ and $\beta = \alpha$ the terms merge into the reduced model
  \[ L(N; \psi, \alpha, E) = \psi N^{-\alpha} + E, \qquad \psi := A + B k^{-\alpha}. \]
  One ray pins down $\psi$, not $A$ and $B$: a ``sloppy'' model; $\varepsilon := |\alpha - \beta|$
  (Chinchilla: $\alpha \approx 0.34$, $\beta \approx 0.28$).
\end{block}

\begin{block}{Gauss-Newton conditioning}
  Minimize $S(\boldsymbol{\theta}) := \tfrac{1}{2} \lVert \mathbf{r}(\boldsymbol{\theta}) \rVert^2$ with Jacobian
  $J_{ij} = \partial r_i / \partial \theta_j$ and steps $(J^\top J)\, \Delta\boldsymbol{\theta} = -J^\top \mathbf{r}$;
  $\kappa_{A,B} := \kappa\bigl((J^\top J)_{A,B}\bigr)$. On a ray the $A$, $B$ columns are, up to sign,
  $N^{-\alpha}$ and $k^{-\beta} N^{-\beta}$. Each law has such a near-parallel pair.
\end{block}

\end{column}
\separatorcolumn

% ==== column 2
\begin{column}{\colwidth}



\begin{block}{One ray is ill-conditioned}
  \statement{Proposition 2 (four laws), Corollary 2 (Chinchilla, Gaussian noise)} On one ray,
  as $\varepsilon \to 0$:
  \[ \kappa_{A,B} = \Theta(\varepsilon^{-2}), \qquad
     \frac{\mathrm{CI}_{0.95}(A)}{\mathrm{CI}_{0.95}(\psi)} = \Omega(\varepsilon^{-1}). \]
  \begin{center}
  \begin{tabular}{@{}lccc@{}}
    \toprule
    Law & $\varepsilon$ & $\varepsilon^{-1}$ & $\varepsilon^{-2}$ \\
    \midrule
    Chinchilla & 0.06 & 17 & 278 \\
    Kaplan & 0.20 & 5 & 25 \\
    Droppo-Elibol & 0.18-0.80 & 1.3-5.6 & 1.6-32 \\
    \bottomrule
  \end{tabular}
  \end{center}
\end{block}

\begin{alertblock}{TPP-diversity threshold}
  \statement{Proposition 1} Chinchilla and unsaturated repeated-data (Kaplan, Droppo-Elibol: a heuristic).
  To leading order, $\kappa_{A,B} \leq \kappa_{\mathrm{target}}$ whenever
  \[ \underbrace{\frac{1}{K} \sum_{\ell=1}^{K} k_\ell^{-2\beta_{\mathrm{eff}}}
     - \Bigl( \frac{1}{K} \sum_{\ell=1}^{K} k_\ell^{-\beta_{\mathrm{eff}}} \Bigr)^{2}}_{V_K}
     \;\geq\;
     \underbrace{\frac{\bigl( K + \sum_{\ell=1}^{K} k_\ell^{-2\beta_{\mathrm{eff}}} \bigr)^{2}}{K^{2}\, \kappa_{\mathrm{target}}}}_{\tau_K}. \]
  Checkable before training. One ray never passes: $V_1 = 0 < \tau_1$.
\end{alertblock}

\begin{block}{The threshold on published designs}
  \begin{center}
  \begin{tabular}{@{}lccc@{}}
    \toprule
    Design & $K$ & $R = k_K / k_1$ & $V_K / \tau_K$ \\
    \midrule
    Hoffmann et al.\ (2022) & 151 & $9.5 \times 10^3$ & \textbf{5.55} \\
    Hu et al.\ (2024) & 6 & 6 & 0.30 \\
    Bhagia et al.\ (2025) & 4 & 10 & 0.47 \\
    \bottomrule
  \end{tabular}
  \end{center}
  At $\kappa_{\mathrm{target}} = 100$ ($\geq 1$ passes), 11 of 24 published designs pass and 13 fail; all 6
  collinear ones fail. Our \CO{} fan: 0.40 at $\beta_{\mathrm{eff}} = 0.28$.
\end{block}

\begin{block}{Holdout prediction}
  \statement{Theorem 1} Any law of Table 2 (repeated-data unsaturated), Gaussian residuals, \CO{} on one
  ray, holdout $H$ with off-ratio points; at leading order in $\sigma$, as $\varepsilon \to 0$:
  \[ \mathbb{E}\bigl[\mathrm{RMSE}^{\mathrm{NC}}_{H}\bigr] < \mathbb{E}\bigl[\mathrm{RMSE}^{\mathrm{CO}}_{H}\bigr]. \]
\end{block}

\begin{block}{Below the threshold, and the stakes}
  Below the threshold, \NC{} beats budget-matched \CO{} in \textbf{68.5\%} (Chinchilla), \textbf{57.4\%}
  (Droppo-Elibol), \textbf{80.6\%} (Kaplan); CIs exclude 50\%. On Wikipedia, \CO{} and \NC{} Chinchilla fits
  imply $N^* \propto C^{0.36}$ and $N^* \propto C^{0.68}$.
\end{block}

\end{column}
\separatorcolumn

% ==== column 3
\begin{column}{\colwidth}

\begin{block}{\NC{} beats \CO{} on held-out runs}
  \begin{columns}[t]
    \begin{column}{0.48\linewidth}
      \centering
      {\fontsize{120}{128}\selectfont\bfseries\textcolor{accent}{97.3\%}}\par\smallskip
      {\large FP32: 1,460/1,500}
    \end{column}
    \begin{column}{0.48\linewidth}
      \centering
      {\fontsize{120}{128}\selectfont\bfseries\textcolor{accent}{98.7\%}}\par\smallskip
      {\large BF16, Wikipedia: 296/300}
    \end{column}
  \end{columns}
  \medskip
  Seed-paired; 4 laws $\times$ 5 corpora, $\sim$1,900 LLaMA-style models of 5.04-76.5M parameters.
  Holdout $R^2$: \CO{} 0.837, \NC{} 0.932; \NC{} wins 98.7\% off the fan ($H_{\mathrm{nc}}$) and 72.1\% along
  it ($H_{\mathrm{col}}$).
\end{block}

\begin{block}{Kaplan law on RedPajama, first epoch}
  \begin{columns}[t]
    \begin{column}{0.49\linewidth}
      \centering
      \includegraphics[width=\linewidth]{01_convergence_Red_Pajama_Kaplan_First.pdf}\\
      {\small Holdout $R^2$ against TPP coverage, 30 seeds: \CO{} blue, \NC{} red.}
    \end{column}
    \begin{column}{0.49\linewidth}
      \centering
      \includegraphics[width=\linewidth]{01_isoflop_Red_Pajama_Kaplan_First.pdf}\\
      {\small Induced isoFLOP curves, seed 0: \CO{} dashed, \NC{} solid, coloured by $C = 6ND$.}
    \end{column}
  \end{columns}
  \medskip
  For Kaplan this error is bias, not conditioning.
\end{block}



\begin{block}{Design recipe}
  \begin{enumerate}
    \item Pick $k_1$ and the minimum spread $R$ at $\kappa_{\mathrm{target}} = 100$.
    \item Train on $K = 2$ ratios at its two ends, runs split evenly, log-uniform in $N$.
    \item Check the exact $\kappa_{A,B}$ before training.
  \end{enumerate}
  20 runs at $k \in \{20, 400\}$: $V_2 / \tau_2 = 1.22$, passes; $\{20, 100\}$: 0.48, fails.
\end{block}

\begin{alertblock}{Takeaways}
  \begin{itemize}
    \item \textbf{A fixed TPP ratio} prescribes how to train one model, not how to estimate a scaling law.
    \item \textbf{On one ray}, report $\psi$, not separate model and data contributions.
  \end{itemize}
  {\small \textbf{Limitations:} 5.04-76.5M parameters, one LLaMA-style architecture, English text,
  pre-training loss; Theorem 1 is proved for one ray, $K \geq 2$ is shown empirically.}
\end{alertblock}

\end{column}
\separatorcolumn
\end{columns}
\end{frame}
\end{document}
